\documentclass[11pt,a4paper]{article}
\usepackage[margin=1in]{geometry}
\usepackage{amsmath,amssymb,amsthm}
\usepackage[hidelinks]{hyperref}

\newtheorem{theorem}{Theorem}
\newtheorem{proposition}[theorem]{Proposition}
\newtheorem{conjecture}[theorem]{Conjecture}

\DeclareMathOperator{\ilen}{\ell_{\mathrm{inv}}}
\DeclareMathOperator{\rank}{rank}

\title{Counterexamples to the Hart--Rowley Involution-Length Conjecture in Finite Coxeter Groups}
\author{Salem Levinstein}
\date{November 27, 2026}

\begin{document}
\maketitle

\begin{abstract}
Hart and Rowley (\texttt{arXiv:2502.00777v1}, \emph{Archiv der Mathematik} 2026) conjectured that, assuming every finite Coxeter group $W$ satisfies the ancestor property (Conjecture~1.2, recently proved in full generality by Gaetz and Gao, \texttt{arXiv:2610.22857}), the greedy ancestor-decomposition involution length satisfies $\ilen(w) \le \rank(W)$ for every $w \in W$ (Conjecture~1.3). We refute Conjecture~1.3 by exhibiting explicit counterexamples in $H_3$ ($\ilen = 4 > 3$), $F_4$ ($5 > 4$), and $E_6$ ($8 > 6$), with parabolic inheritance ($E_6 \subset E_7$) immediately extending the failure to $E_7$ ($\max \ilen \ge 8 > 7$). In particular, the computational verification of Conjecture~1.3 reported for $F_4, E_6, H_3$ in Section~3 of \texttt{arXiv:2502.00777v1} is incorrect under the paper's definitions. The complete $H_3$ refutation---including unconditional verification of the ancestor property on all $119$ non-identity elements of any \texttt{CoxeterSystem CoxeterMatrix.H$_3$}---is formally machine-checked in Lean~4 (\texttt{v4.33.1} + Mathlib, standard axioms only).
\end{abstract}

\section{Introduction}

Let $(W, R)$ be a Coxeter system with Coxeter length function $\ell : W \to \mathbb{N}$ and rank $\rank(W) = |R|$. Following Hart and Rowley \cite{HartRowley2025}, an element $u \in W$ is a \emph{prefix} of $w \in W$ if $\ell(u) + \ell(u^{-1}w) = \ell(w)$, equivalently if $u$ lies in the right weak order interval $[1, w]_R$. Let
\[
\mathrm{ipref}(w) = \{u \in W : u \text{ is a prefix of } w \text{ and } o(u) = 2\}
\]
denote the set of involution prefixes of $w$, and let $A(w) \subseteq \mathrm{ipref}(w)$ be the subset of elements of maximal Coxeter length. The group $W$ has the \emph{ancestor property} if $|A(w)| = 1$ for every $w \in W \setminus \{1\}$, in which case the unique element $a(w) \in A(w)$ is called the \emph{ancestor} of $w$.

Hart and Rowley conjectured that every finite Coxeter group has the ancestor property (Conjecture~1.2 of \cite{HartRowley2025}); this conjecture has recently been proved for all Coxeter groups by Gaetz and Gao \cite{GaetzGao2026}. Whenever $W$ has the ancestor property, every element $w \in W$ admits a unique greedy \emph{ancestor decomposition} $w = a_1 a_2 \cdots a_k$ defined recursively by $a_1 = a(w)$ and applying the decomposition to $a_1^{-1} w = a_1 w$, with $\ilen(1) = 0$ and $\ilen(w) = k$ called the \emph{involution length} of $w$.

\begin{conjecture}[Hart--Rowley \cite{HartRowley2025}, Conjecture~1.3]\label{conj:hr}
Suppose Conjecture~1.2 holds, and let $W$ be a finite Coxeter group. For each $w \in W$, $\ilen(w) \le \rank(W)$.
\end{conjecture}

In Section~3 of \cite{HartRowley2025}, the authors reported that Conjectures~1.2 and~1.3 had been verified in Magma for all irreducible finite Coxeter groups of order less than $100{,}000$, including $H_3, F_4, H_4, E_6$. We show that while Conjecture~1.2 indeed holds on $H_3, F_4, E_6$, Conjecture~1.3 is false already in $H_3$, $F_4$, and $E_6$ (and hence also in $E_7$).

\section{Main Counterexample in $H_3$}

Let $W = H_3$ be the finite Coxeter group of rank $3$ and order $120$ generated by $R = \{r_1, r_2, r_3\}$ with Coxeter matrix $m(r_1, r_2) = 3$, $m(r_2, r_3) = 5$, and $m(r_1, r_3) = 2$.

\begin{theorem}\label{thm:h3}
The finite Coxeter group $H_3$ satisfies the ancestor property, and the length-$7$ element
\[
w_1 = r_2 r_3 r_1 r_2 r_3 r_2 r_1 \in H_3
\]
has exact ancestor decomposition
\[
w_1 = (r_2 r_3 r_1 r_2)(r_3)(r_2)(r_1), \qquad \ilen(w_1) = 4 > 3 = \rank(H_3).
\]
In particular, Conjecture~\ref{conj:hr} is false.
\end{theorem}

\begin{proof}
Because $r_1$ and $r_3$ commute, $a_1 := r_2 r_3 r_1 r_2$ satisfies $a_1^{-1} = r_2 r_1 r_3 r_2 = r_2 r_3 r_1 r_2 = a_1$, and $a_1 \neq 1$ has Coxeter length $4$. Direct enumeration of the right weak order interval $[1, w_1]_R$ (verified both in Lean~4 and in Python) shows that $w_1$ has $\ell(w_1) = 7$ and exactly $9$ prefixes:
\[
1,\quad r_2,\quad r_2 r_1,\quad r_2 r_3,\quad r_2 r_3 r_1,\quad r_2 r_3 r_1 r_2,\quad r_2 r_3 r_1 r_2 r_3,\quad r_2 r_3 r_1 r_2 r_3 r_2,\quad w_1.
\]
Of these $9$ prefixes, only $r_2$ (length $1$) and $a_1 = r_2 r_3 r_1 r_2$ (length $4$) are involutions. Hence $A(w_1) = \{a_1\}$ and $a(w_1) = a_1$.

Stripping off $a_1$ leaves the Coxeter element $c := a_1^{-1} w_1 = r_3 r_2 r_1$ of length $3$. Its prefixes are $1, r_3, r_3 r_2, r_3 r_2 r_1$. Since $o(r_3 r_2) = 5$ and $o(r_3 r_2 r_1) = 10$, neither is an involution, so $a(c) = r_3$. Removing $r_3$ leaves $r_2 r_1$, whose prefixes are $1, r_2, r_2 r_1$; since $o(r_2 r_1) = 3$, its ancestor is $r_2$. Finally $r_2^{-1}(r_2 r_1) = r_1$ is an involution, so $a(r_1) = r_1$.

Hence $w_1 = (r_2 r_3 r_1 r_2)(r_3)(r_2)(r_1)$ is the unique ancestor decomposition of $w_1$, so $\ilen(w_1) = 4 > 3 = \rank(H_3)$.
\end{proof}

A second independent length-$7$ witness in $H_3$ is
\[
w_2 = r_2 r_3 r_2 r_3 r_1 r_2 r_3 = (r_2 r_3 r_2)(r_3 r_1)(r_2)(r_3), \qquad \ilen(w_2) = 4 > 3.
\]
Exhaustive checking over all $119$ non-identity elements of $H_3$ (together with $\ilen(1) = 0$) confirms $\max_{w \in H_3} \ilen(w) = 4$.

\section{Crystallographic Counterexamples and Parabolic Inheritance}

For any standard parabolic subgroup $W_J = \langle J \rangle \le W$ ($J \subseteq R$), every reduced expression of $w \in W_J$ uses only letters in $J$, so every prefix of $w$ in $W$ belongs to $W_J$ with the same Coxeter length. Consequently, $\ell_{\mathrm{inv}, W}(w) = \ell_{\mathrm{inv}, W_J}(w)$ for all $w \in W_J$.

Exact root-system computation over $\mathbb{Z}$ (\texttt{verify.py}) confirms Conjecture~1.2 on all elements of $F_4$ and $E_6$ and produces explicit crystallographic counterexamples to Conjecture~\ref{conj:hr}:
\begin{itemize}
\item \textbf{Type $F_4$} (rank $4$, order $1{,}152$, Bourbaki diagram $1 - 2 \Rightarrow 3 - 4$):
\[
w = r_3 r_4 r_1 r_2 r_3 r_2 r_1 = (r_3 r_1)(r_4 r_2)(r_3)(r_2)(r_1), \qquad \ilen(w) = 5 > 4,
\]
as well as $r_4 r_2 r_3 r_1 r_2 r_3 r_4 = (r_4 r_2)(r_3 r_1)(r_2)(r_3)(r_4)$, with $\max_{w \in F_4} \ilen(w) = 5$.
\item \textbf{Type $E_6$} (rank $6$, order $51{,}840$, Bourbaki labeling with $2$ adjacent to $4$):
\[
w = r_3 r_2 r_4 r_5 r_6 r_1 r_3 r_4 r_5 r_2 r_4 r_3 r_1 = (r_3 r_2)(r_4 r_1)(r_5 r_3)(r_6 r_4)(r_5 r_2)(r_4)(r_3)(r_1),
\]
giving $\ilen(w) = 8 > 6$. Exhaustive search gives $\max_{w \in E_6} \ilen(w) = 8$.
\end{itemize}
By parabolic inheritance ($E_6 \subset E_7$), Conjecture~\ref{conj:hr} also fails in $E_7$, with $\max_{w \in E_7} \ilen(w) \ge 8 > 7 = \rank(E_7)$.

\section{Formal Verification in Lean~4}

Theorem~\ref{thm:h3} (along with the second witness $w_2$ and an unconditional verification of Conjecture~1.2 on all $119$ non-identity elements of $H_3$) is formally proved in Lean~4 (\texttt{v4.33.1} + Mathlib) in the file \texttt{HartRowleyCoxeterInvolutionLength.lean} for every abstract \texttt{CoxeterSystem CoxeterMatrix.H$_3$ W}, using only the three standard axioms (\texttt{propext}, \texttt{Classical.choice}, \texttt{Quot.sound}), with zero \texttt{sorry}, \texttt{admit}, or \texttt{native\_decide}.

Specifically, we lift \texttt{CoxeterMatrix.H$_3$} to the $30$-root permutation representation $\varphi : W \to \mathrm{Perm}(\mathrm{Fin}\,30)$, prove $\mathrm{invCount}(\varphi(w)) \le \ell(w)$ from the single-root flip property of simple reflections, certify right multiplication on $120$ canonical reduced words via a reflection-verified braid/cancellation rewriter (giving surjectivity of $\mathrm{Fin}\,120 \to W$), deduce injectivity from distinctness of the action on the three simple roots, and thereby obtain exact Coxeter lengths, the ancestor property on all of $H_3$, and $\ilen(w_1) = \ilen(w_2) = 4 > 3$.

\begin{thebibliography}{2}
\bibitem{HartRowley2025}
S. B. Hart and P. J. Rowley,
\emph{A note on involution prefixes in Coxeter groups},
\emph{Archiv der Mathematik} (2026), arXiv:2502.00777v1 [math.GR], doi:10.1007/s00013-025-02220-z.

\bibitem{GaetzGao2026}
C. Gaetz and Y. Gao,
\emph{Separable elements and Bruhat orders},
arXiv:2610.22857 [math.CO], October 2026.
\end{thebibliography}

\end{document}
